\DocumentMetadata{
  pdfstandard={ua-2},
  tagging=on,
  lang=en,
  xmp=true}
% Options for packages loaded elsewhere
% Options for packages loaded elsewhere
\PassOptionsToPackage{unicode}{hyperref}
\PassOptionsToPackage{hyphens}{url}
\PassOptionsToPackage{dvipsnames,svgnames,x11names}{xcolor}
%
\DocumentMetadata{
  pdfstandard=ua-2,
  lang=en-US,
  tagging=on,
  tagging-setup={math/alt/use=true}
}
% Options for packages loaded elsewhere
\PassOptionsToPackage{unicode}{hyperref}
\PassOptionsToPackage{hyphens}{url}
\PassOptionsToPackage{dvipsnames,svgnames,x11names}{xcolor}
\documentclass[
  english,
]{article}
\usepackage{xcolor}
\usepackage{amsmath,amssymb}
\setcounter{secnumdepth}{-\maxdimen} % remove section numbering
\usepackage{iftex}
\ifPDFTeX
  \usepackage[T1]{fontenc}
  \usepackage[utf8]{inputenc}
  \usepackage{textcomp} % provide euro and other symbols
\else % if luatex or xetex
  \usepackage{unicode-math} % this also loads fontspec
  \defaultfontfeatures{Scale=MatchLowercase}
  \defaultfontfeatures[\rmfamily]{Ligatures=TeX,Scale=1}
\fi
\usepackage{lmodern}
\ifPDFTeX\else
  % xetex/luatex font selection
  \setmainfont[]{Libertinus Serif}
  \setsansfont[]{Libertinus Sans}
  \setmonofont[Scale=0.95]{Latin Modern Mono}
  \setmathfont[]{Libertinus Math}
\fi
% Use upquote if available, for straight quotes in verbatim environments
\IfFileExists{upquote.sty}{\usepackage{upquote}}{}
\IfFileExists{microtype.sty}{% use microtype if available
  \usepackage[]{microtype}
  \UseMicrotypeSet[protrusion]{basicmath} % disable protrusion for tt fonts
}{}
\makeatletter
\@ifundefined{KOMAClassName}{% if non-KOMA class
  \IfFileExists{parskip.sty}{%
    \usepackage{parskip}
  }{% else
    \setlength{\parindent}{0pt}
    \setlength{\parskip}{6pt plus 2pt minus 1pt}}
}{% if KOMA class
  \KOMAoptions{parskip=half}}
\makeatother
% Make \paragraph and \subparagraph free-standing
\makeatletter
\ifx\paragraph\undefined\else
  \let\oldparagraph\paragraph
  \renewcommand{\paragraph}{
    \@ifstar
      \xxxParagraphStar
      \xxxParagraphNoStar
  }
  \newcommand{\xxxParagraphStar}[1]{\oldparagraph*{#1}\mbox{}}
  \newcommand{\xxxParagraphNoStar}[1]{\oldparagraph{#1}\mbox{}}
\fi
\ifx\subparagraph\undefined\else
  \let\oldsubparagraph\subparagraph
  \renewcommand{\subparagraph}{
    \@ifstar
      \xxxSubParagraphStar
      \xxxSubParagraphNoStar
  }
  \newcommand{\xxxSubParagraphStar}[1]{\oldsubparagraph*{#1}\mbox{}}
  \newcommand{\xxxSubParagraphNoStar}[1]{\oldsubparagraph{#1}\mbox{}}
\fi
\makeatother

\usepackage{color}
\usepackage{fancyvrb}
\newcommand{\VerbBar}{|}
\newcommand{\VERB}{\Verb[commandchars=\\\{\}]}
\DefineVerbatimEnvironment{Highlighting}{Verbatim}{commandchars=\\\{\}}
% Add ',fontsize=\small' for more characters per line
\usepackage{framed}
\definecolor{shadecolor}{RGB}{241,243,245}
\newenvironment{Shaded}{\begin{snugshade}}{\end{snugshade}}
\newcommand{\AlertTok}[1]{\textcolor[rgb]{0.68,0.00,0.00}{#1}}
\newcommand{\AnnotationTok}[1]{\textcolor[rgb]{0.37,0.37,0.37}{#1}}
\newcommand{\AttributeTok}[1]{\textcolor[rgb]{0.40,0.45,0.13}{#1}}
\newcommand{\BaseNTok}[1]{\textcolor[rgb]{0.68,0.00,0.00}{#1}}
\newcommand{\BuiltInTok}[1]{\textcolor[rgb]{0.00,0.23,0.31}{#1}}
\newcommand{\CharTok}[1]{\textcolor[rgb]{0.13,0.47,0.30}{#1}}
\newcommand{\CommentTok}[1]{\textcolor[rgb]{0.37,0.37,0.37}{#1}}
\newcommand{\CommentVarTok}[1]{\textcolor[rgb]{0.37,0.37,0.37}{\textit{#1}}}
\newcommand{\ConstantTok}[1]{\textcolor[rgb]{0.56,0.35,0.01}{#1}}
\newcommand{\ControlFlowTok}[1]{\textcolor[rgb]{0.00,0.23,0.31}{\textbf{#1}}}
\newcommand{\DataTypeTok}[1]{\textcolor[rgb]{0.68,0.00,0.00}{#1}}
\newcommand{\DecValTok}[1]{\textcolor[rgb]{0.68,0.00,0.00}{#1}}
\newcommand{\DocumentationTok}[1]{\textcolor[rgb]{0.37,0.37,0.37}{\textit{#1}}}
\newcommand{\ErrorTok}[1]{\textcolor[rgb]{0.68,0.00,0.00}{#1}}
\newcommand{\ExtensionTok}[1]{\textcolor[rgb]{0.00,0.23,0.31}{#1}}
\newcommand{\FloatTok}[1]{\textcolor[rgb]{0.68,0.00,0.00}{#1}}
\newcommand{\FunctionTok}[1]{\textcolor[rgb]{0.28,0.35,0.67}{#1}}
\newcommand{\ImportTok}[1]{\textcolor[rgb]{0.00,0.46,0.62}{#1}}
\newcommand{\InformationTok}[1]{\textcolor[rgb]{0.37,0.37,0.37}{#1}}
\newcommand{\KeywordTok}[1]{\textcolor[rgb]{0.00,0.23,0.31}{\textbf{#1}}}
\newcommand{\NormalTok}[1]{\textcolor[rgb]{0.00,0.23,0.31}{#1}}
\newcommand{\OperatorTok}[1]{\textcolor[rgb]{0.37,0.37,0.37}{#1}}
\newcommand{\OtherTok}[1]{\textcolor[rgb]{0.00,0.23,0.31}{#1}}
\newcommand{\PreprocessorTok}[1]{\textcolor[rgb]{0.68,0.00,0.00}{#1}}
\newcommand{\RegionMarkerTok}[1]{\textcolor[rgb]{0.00,0.23,0.31}{#1}}
\newcommand{\SpecialCharTok}[1]{\textcolor[rgb]{0.37,0.37,0.37}{#1}}
\newcommand{\SpecialStringTok}[1]{\textcolor[rgb]{0.13,0.47,0.30}{#1}}
\newcommand{\StringTok}[1]{\textcolor[rgb]{0.13,0.47,0.30}{#1}}
\newcommand{\VariableTok}[1]{\textcolor[rgb]{0.07,0.07,0.07}{#1}}
\newcommand{\VerbatimStringTok}[1]{\textcolor[rgb]{0.13,0.47,0.30}{#1}}
\newcommand{\WarningTok}[1]{\textcolor[rgb]{0.37,0.37,0.37}{\textit{#1}}}

\usepackage{longtable,booktabs,array}
\newcounter{none} % for unnumbered tables
\usepackage{calc} % for calculating minipage widths
% Correct order of tables after \paragraph or \subparagraph
\usepackage{etoolbox}
\makeatletter
\patchcmd\longtable{\par}{\if@noskipsec\mbox{}\fi\par}{}{}
\makeatother
% Allow footnotes in longtable head/foot
\IfFileExists{footnotehyper.sty}{\usepackage{footnotehyper}}{\usepackage{footnote}}
\makesavenoteenv{longtable}
\usepackage{graphicx}
\makeatletter
\newsavebox\pandoc@box
\newcommand*\pandocbounded[1]{% scales image to fit in text height/width
  \sbox\pandoc@box{#1}%
  \Gscale@div\@tempa{\textheight}{\dimexpr\ht\pandoc@box+\dp\pandoc@box\relax}%
  \Gscale@div\@tempb{\linewidth}{\wd\pandoc@box}%
  \ifdim\@tempb\p@<\@tempa\p@\let\@tempa\@tempb\fi% select the smaller of both
  \ifdim\@tempa\p@<\p@\scalebox{\@tempa}{\usebox\pandoc@box}%
  \else\usebox{\pandoc@box}%
  \fi%
}
% Set default figure placement to htbp
\def\fps@figure{htbp}
\makeatother



\ifLuaTeX
\usepackage[bidi=basic,shorthands=off]{babel}
\else
\usepackage[bidi=default,shorthands=off]{babel}
\fi
\ifPDFTeX
\else
\babelfont{rm}[]{Libertinus Serif}
\fi
\ifLuaTeX
  \usepackage{selnolig} % disable illegal ligatures
\fi


\setlength{\emergencystretch}{3em} % prevent overfull lines

\providecommand{\tightlist}{%
  \setlength{\itemsep}{0pt}\setlength{\parskip}{0pt}}



 




\usepackage{physics}
\usepackage{braket}

\newcommand{\dl}[1]{{\hspace{#1mu}\mathrm d}}
\newcommand{\me}{{\mathrm e}}

\newcommand{\rss}{\mathrm{RSS}}
\newcommand{\ols}{\mathrm{OLS}}
\newcommand{\ridge}{\mathrm{ridge}}
\newcommand{\lasso}{\mathrm{lasso}}
\newcommand{\mle}{\mathrm{MLE}}
\newcommand{\constant}{\mathrm{constant}}
\newcommand{\resid}{\mathrm{residual}}

\newcommand{\Exp}{\operatorname{E}}
\newcommand{\Var}{\operatorname{Var}}
\newcommand{\Mode}{\operatorname{mode}}
\newcommand{\cor}{\operatorname{cor}}
\newcommand{\nll}{\operatorname{nll}}
\newcommand{\span}{\operatorname{span}}

\newcommand{\pdfbinom}{{\tt binom}}
\newcommand{\pdfbeta}{{\tt beta}}
\newcommand{\pdfpois}{{\tt poisson}}
\newcommand{\pdfgamma}{{\tt gamma}}
\newcommand{\pdfnormal}{{\tt norm}}
\newcommand{\pdfexp}{{\tt expon}}

\newcommand{\distbinom}{\operatorname{B}}
\newcommand{\distbeta}{\operatorname{Beta}}
\newcommand{\distgamma}{\operatorname{Gamma}}
\newcommand{\distexp}{\operatorname{Exp}}
\newcommand{\distpois}{\operatorname{Poisson}}
\newcommand{\distnormal}{\operatorname{\mathcal N}}


\usepackage{xcolor}
\definecolor{AccessibleRed}{HTML}{BC4B00}
\newcommand{\red}[1]{\textcolor{AccessibleRed}{\textbf{#1}}}
\usepackage{alltt}
\usepackage{enumitem}
\usepackage[letterpaper,top=1in,bottom=0.89in,left=1in,right=1in]{geometry}
\AtBeginDocument{
    \ifdefined\theexercise
        \renewcommand{\theexercise}{\arabic{exercise}}
    \fi
}           
\newcommand{\mc}[1]{
    \begin{enumerate}
        \setcounter{enumi}{0}
        \renewcommand{\labelenumi}{\Alph{enumi}.}
        \setlength{\itemsep}{0pt}
        \setlength{\topsep}{2pt}
        #1
    \end{enumerate}
}
\usepackage{multicol}
\newcommand{\mcc}[2][2]{
    \begin{multicols}{#1}
        \begin{enumerate}
            \setcounter{enumi}{0}
            \renewcommand{\labelenumi}{\Alph{enumi}.}
            \setlength{\itemsep}{0pt}
            \setlength{\topsep}{2pt}
            #2
        \end{enumerate}
    \end{multicols}
}
\makeatletter
\@ifpackageloaded{caption}{}{\usepackage{caption}}
\AtBeginDocument{%
\ifdefined\contentsname
  \renewcommand*\contentsname{Table of contents}
\else
  \newcommand\contentsname{Table of contents}
\fi
\ifdefined\listfigurename
  \renewcommand*\listfigurename{List of Figures}
\else
  \newcommand\listfigurename{List of Figures}
\fi
\ifdefined\listtablename
  \renewcommand*\listtablename{List of Tables}
\else
  \newcommand\listtablename{List of Tables}
\fi
\ifdefined\figurename
  \renewcommand*\figurename{Figure}
\else
  \newcommand\figurename{Figure}
\fi
\ifdefined\tablename
  \renewcommand*\tablename{Table}
\else
  \newcommand\tablename{Table}
\fi
}
\@ifpackageloaded{float}{}{\usepackage{float}}
\floatstyle{ruled}
\@ifundefined{c@chapter}{\newfloat{codelisting}{h}{lop}}{\newfloat{codelisting}{h}{lop}[chapter]}
\floatname{codelisting}{Listing}
\newcommand*\listoflistings{\listof{codelisting}{List of Listings}}
\usepackage{amsthm}
\theoremstyle{definition}
\newtheorem{example}{Example}[section]
\theoremstyle{definition}
\newtheorem{exercise}{Exercise}[section]
\theoremstyle{remark}
\AtBeginDocument{\renewcommand*{\proofname}{Proof}}
\newtheorem*{remark}{Remark}
\newtheorem*{solution}{Solution}
\newtheorem{refremark}{Remark}[section]
\newtheorem{refsolution}{Solution}[section]
\makeatother
\makeatletter
\makeatother
\makeatletter
\@ifpackageloaded{caption}{}{\usepackage{caption}}
\@ifpackageloaded{subcaption}{}{\usepackage{subcaption}}
\makeatother
\usepackage{bookmark}
\IfFileExists{xurl.sty}{\usepackage{xurl}}{} % add URL line breaks if available
\urlstyle{same}
\hypersetup{
  pdftitle={Homework 2},
  pdflang={en},
  colorlinks=true,
  linkcolor={blue},
  filecolor={Maroon},
  citecolor={Blue},
  urlcolor={Blue},
  pdfcreator={LaTeX via pandoc}}


\title{Homework 2}
\author{}
\date{}
\begin{document}
\maketitle


\section{Week 2: Regularization
Methods}\label{week-2-regularization-methods}

\subsection{Theory Problems}\label{theory-problems}

\begin{exercise}[]\protect\hypertarget{exr-}{}\label{exr-}

Suppose we estimate the regression coefficients in a linear regression
model by minimizing \[
\sum_{i=1}^n\qty(y_i-\beta_0-\sum_{j=1}^p\beta_jx_{ij})^2+\lambda\sum_{j=1}^p\beta_j^2
\] for a particular value of \(\lambda\). Answer the following questions
and justify your answer.

\begin{enumerate}
\def\labelenumi{\alph{enumi}.}
\tightlist
\item
  As we increase \(\lambda\) from \(0\), how will the training \(\rss\)
  be changed?
\item
  As we increase \(\lambda\) from \(0\), how will the test \(\rss\) be
  changed?
\item
  As we increase \(\lambda\) from \(0\), how will the variance be
  changed?
\item
  As we increase \(\lambda\) from \(0\), how will the squared bias be
  changed?
\item
  As we increase \(\lambda\) from \(0\), how will the irreducible error
  be changed?
\end{enumerate}

\end{exercise}

\begin{solution}
This is ridge regression. The tuning parameter \(\lambda\) controls the
strength of the penalty:

\[
\lambda \sum_{j=1}^p \beta_j^2.
\]

When \(\lambda=0\), the method is ordinary least squares. As \(\lambda\)
increases, the penalty becomes stronger, so the coefficients are shrunk
more toward zero.

\textbf{a. Training RSS}

As \(\lambda\) increases from \(0\), the training RSS will increase or
stay the same.

Reason: \(\lambda=0\) gives the ordinary least squares solution, which
minimizes the training RSS without penalty. Adding a penalty restricts
the solution away from the unconstrained RSS minimizer, so the training
RSS cannot improve.

\textbf{b. Test RSS}

As \(\lambda\) increases from \(0\), the test RSS will typically
decrease at first and then increase.

Reason: a small amount of shrinkage can reduce overfitting and lower
variance, improving test performance. But if \(\lambda\) becomes too
large, the model becomes too constrained, leading to underfitting and
larger bias.

\textbf{c.~Variance}

As \(\lambda\) increases from \(0\), variance will decrease.

Reason: larger \(\lambda\) shrinks the coefficients more strongly toward
zero, making the fitted model less sensitive to random variation in the
training data.

\textbf{d.~Squared bias}

As \(\lambda\) increases from \(0\), squared bias will increase.

Reason: stronger shrinkage pulls the fitted model away from the least
squares fit and may prevent it from capturing the true relationship
fully.

\textbf{e. Irreducible error}

The irreducible error does not change.

Reason: irreducible error is caused by the noise in the data-generating
process, not by the choice of \(\lambda\) or the regression method.
\end{solution}

\begin{exercise}[]\protect\hypertarget{exr-}{}\label{exr-}

Suppose we estimate the regression coefficients in a linear regression
model by minimizing \[
\sum_{i=1}^n\qty(y_i-\beta_0-\sum_{j=1}^p\beta_jx_{ij})^2\quad \text{ subject to }\sum_{j=1}^p\abs{\beta_j}\leq s
\] for a particular value of \(s\). Answer the following questions and
justify your answer.

\begin{enumerate}
\def\labelenumi{\alph{enumi}.}
\tightlist
\item
  As we increase \(s\) from \(0\), how will the training \(\rss\) be
  changed?
\item
  As we increase \(s\) from \(0\), how will the test \(\rss\) be
  changed?
\item
  As we increase \(s\) from \(0\), how will the variance be changed?
\item
  As we increase \(s\) from \(0\), how will the squared bias be changed?
\item
  As we increase \(s\) from \(0\), how will the irreducible error be
  changed?
\end{enumerate}

\end{exercise}

\begin{solution}
This is the lasso constraint form. The parameter \(s\) controls how
large the coefficients are allowed to be:

\[
\sum_{j=1}^p |\beta_j| \le s.
\]

When \(s=0\), all slope coefficients must be zero, so the model is the
intercept-only model. As \(s\) increases, the constraint becomes less
restrictive, and the model becomes more flexible.

\textbf{a. Training RSS}

As \(s\) increases, the training RSS will decrease or stay the same.

Reason: increasing \(s\) enlarges the set of allowable coefficient
values. Since the old solution is still feasible, the minimized training
RSS cannot increase.

\textbf{b. Test RSS}

As \(s\) increases, the test RSS will typically decrease at first and
then increase.

Reason: when \(s\) is very small, the model is too simple and has high
bias, so test error is large. Increasing \(s\) allows the model to
capture more signal, reducing test error. However, if \(s\) becomes too
large, the model may overfit the training data, increasing variance and
therefore increasing test error.

\textbf{c.~Variance}

As \(s\) increases, variance will increase.

Reason: larger \(s\) allows more flexible models with larger
coefficients. The fitted model becomes more sensitive to the particular
training sample, so its variance increases.

\textbf{d.~Squared bias}

As \(s\) increases, squared bias will decrease.

Reason: when \(s\) is small, the model is heavily constrained and may
underfit the true relationship. Increasing \(s\) relaxes the constraint,
allowing the model to better approximate the true regression function.

\textbf{e. Irreducible error}

The irreducible error does not change.

Reason: irreducible error comes from the random noise in the
data-generating process, not from the choice of model or the value of
\(s\).
\end{solution}

\begin{example}[PCR vs PLS]\protect\hypertarget{exm-}{}\label{exm-}

Consider a regression problem with centered data matrix
\(X \in \mathbb{R}^{n \times p}\) and response \(y \in \mathbb{R}^n\).

\textbf{a. Bias--Variance Tradeoff}

Discuss how PCR and PLS differ in terms of bias and variance.

\begin{itemize}
\tightlist
\item
  When might PCR have higher bias than PLS?
\item
  When might PLS have higher variance than PCR?
\end{itemize}

\textbf{b. When PCR Fails}

Construct (or describe) a situation where PCR performs way worse than
PLS. Explain why PCR may perform poorly in this setting.

\textbf{c.~Code examples}

Write Python code to generate the datasets to implement the above two
cases. A dataset is a pair \((X,y)\) with both are \texttt{numpy.array}.

Ensure there is a performance difference. You can get a bonus if there
is at least 1\% difference in \(R^2\). You may fix random seeds or
adjust noise levels to achieve this.

You code should look like this.

\begin{Shaded}
\begin{Highlighting}[]
\ImportTok{import}\NormalTok{ numpy }\ImportTok{as}\NormalTok{ np}

\NormalTok{rng }\OperatorTok{=}\NormalTok{ np.random.default\_rng(}\DecValTok{0}\NormalTok{)}

\CommentTok{\#\# Write your data generating code here to replace the code below}
\NormalTok{n, p }\OperatorTok{=} \DecValTok{100}\NormalTok{, }\DecValTok{100}
\NormalTok{X }\OperatorTok{=}\NormalTok{ rng.normal(size}\OperatorTok{=}\NormalTok{(n, p))}
\NormalTok{y }\OperatorTok{=}\NormalTok{ rng.normal(size}\OperatorTok{=}\NormalTok{n)}
\end{Highlighting}
\end{Shaded}

I will test the dataset using the following setting:

\begin{enumerate}
\def\labelenumi{\arabic{enumi}.}
\tightlist
\item
  Split the dataset into trainning and test set with
  \texttt{test\_size=0.3}, \texttt{random\_state=0}.
\item
  Use \texttt{cv=5} and \texttt{neg\_mean\_squared\_error} to pick
  \texttt{n\_components}. You may set the maximal search cap to be
  \texttt{20} for simplicity.
\item
  Train the model on the training set, and test the model on the test
  set.
\item
  For performance I only look at \texttt{r2\_score}.
\end{enumerate}

\end{example}

\begin{solution}
\textbf{(a) Bias--Variance Tradeoff}

PCR and PLS differ in how they construct components, which leads to
different bias--variance behavior.

\begin{itemize}
\tightlist
\item
  \textbf{PCR} selects components based only on the variance of \(X\),
  ignoring \(y\).
\item
  \textbf{PLS} selects components based on their covariance with \(y\).
\end{itemize}

\textbf{Bias}

\begin{itemize}
\item
  PCR may have \textbf{higher bias} when the directions that are most
  predictive of \(y\) correspond to \textbf{low-variance directions in
  \(X\)}.\\
  In this case, PCR may discard these directions when truncating
  components.
\item
  PLS tends to have \textbf{lower bias for prediction}, since it uses
  \(y\) to identify directions that are directly relevant for predicting
  \(y\).
\end{itemize}

\textbf{Variance}

\begin{itemize}
\item
  PLS may have \textbf{higher variance}, because it uses \(y\) to
  construct components.\\
  This makes it more sensitive to noise in \(y\), especially in small
  samples or low signal-to-noise settings.
\item
  PCR often has \textbf{lower variance}, since it relies only on the
  covariance structure of \(X\), which is typically more stable than
  \(X^\top y\).
\end{itemize}

\textbf{Summary:}

\begin{itemize}
\tightlist
\item
  PCR: potentially higher bias, lower variance\\
\item
  PLS: potentially lower bias, higher variance
\end{itemize}

\textbf{(b) When PCR Fails}

PCR performs poorly when the directions with large variance in \(X\) are
not predictive of \(y\), while the predictive signal lies in
low-variance directions.

A typical construction is:

\begin{itemize}
\tightlist
\item
  \(X\) contains:

  \begin{itemize}
  \tightlist
  \item
    high-variance noise directions
  \item
    low-variance signal directions
  \end{itemize}
\item
  \(y\) depends primarily on the low-variance signal
\end{itemize}

In this case:

\begin{itemize}
\tightlist
\item
  PCA identifies the high-variance directions (noise-dominated)
\item
  PCR keeps these directions and may discard the low-variance signal
\item
  As a result, PCR fails to capture the relationship between \(X\) and
  \(y\)
\end{itemize}

\textbf{(c) code example}

\begin{Shaded}
\begin{Highlighting}[]
\ImportTok{import}\NormalTok{ numpy }\ImportTok{as}\NormalTok{ np}

\NormalTok{rng }\OperatorTok{=}\NormalTok{ np.random.default\_rng(}\DecValTok{0}\NormalTok{)}

\NormalTok{n, p }\OperatorTok{=} \DecValTok{200}\NormalTok{, }\DecValTok{100}

\NormalTok{Z\_noise }\OperatorTok{=}\NormalTok{ rng.normal(size}\OperatorTok{=}\NormalTok{(n, }\DecValTok{5}\NormalTok{))}
\NormalTok{A\_noise }\OperatorTok{=}\NormalTok{ rng.normal(size}\OperatorTok{=}\NormalTok{(}\DecValTok{5}\NormalTok{, p))}
\NormalTok{X\_noise }\OperatorTok{=}\NormalTok{ Z\_noise }\OperatorTok{@}\NormalTok{ A\_noise }\OperatorTok{*} \DecValTok{10}

\NormalTok{z\_signal }\OperatorTok{=}\NormalTok{ rng.normal(size}\OperatorTok{=}\NormalTok{(n, }\DecValTok{1}\NormalTok{))}
\NormalTok{a\_signal }\OperatorTok{=}\NormalTok{ rng.normal(size}\OperatorTok{=}\NormalTok{(}\DecValTok{1}\NormalTok{, p))}
\NormalTok{X\_signal }\OperatorTok{=}\NormalTok{ z\_signal }\OperatorTok{@}\NormalTok{ a\_signal }\OperatorTok{*} \FloatTok{0.2}

\NormalTok{X }\OperatorTok{=}\NormalTok{ X\_noise }\OperatorTok{+}\NormalTok{ X\_signal }\OperatorTok{+}\NormalTok{rng.normal(size}\OperatorTok{=}\NormalTok{(n, p))}
\NormalTok{y }\OperatorTok{=}\NormalTok{ z\_signal.ravel() }\OperatorTok{+} \FloatTok{0.5}\OperatorTok{*}\NormalTok{rng.normal(size}\OperatorTok{=}\NormalTok{n)}
\end{Highlighting}
\end{Shaded}

\end{solution}

\begin{example}[Exercise: Ridge vs
Lasso]\protect\hypertarget{exm-}{}\label{exm-}

Consider a regression problem with centered data matrix
\(X \in \mathbb{R}^{n \times p}\) and response \(y \in \mathbb{R}^n\).

\begin{center}\rule{0.5\linewidth}{0.5pt}\end{center}

\subsubsection{(a) Bias--Variance
Tradeoff}\label{a-biasvariance-tradeoff}

Discuss how ridge regression and lasso differ in terms of bias and
variance.

\begin{itemize}
\tightlist
\item
  When might lasso have higher bias than ridge?
\item
  When might ridge have higher variance than lasso?
\end{itemize}

\begin{center}\rule{0.5\linewidth}{0.5pt}\end{center}

\subsubsection{(b) When Ridge Fails}\label{b-when-ridge-fails}

Construct (or describe) a situation where ridge regression performs much
worse than lasso.

Explain your reasoning.

Hint: Consider the case where only a small number of predictors are
truly relevant.

\begin{center}\rule{0.5\linewidth}{0.5pt}\end{center}

\subsubsection{(c) Code Examples}\label{c-code-examples}

Write Python code to generate datasets \((X, y)\) illustrating:

\begin{itemize}
\tightlist
\item
  a case where ridge performs much worse than lasso
\item
  a case where lasso performs worse than ridge
\end{itemize}

Use cross-validation to tune the regularization parameter.

Ensure there is a performance difference. You can get a bonus if there
is at least 1\% difference in \(R^2\). You may fix random seeds or
adjust noise levels to achieve this.

Your code should look like this:

\begin{Shaded}
\begin{Highlighting}[]
\ImportTok{import}\NormalTok{ numpy }\ImportTok{as}\NormalTok{ np}

\NormalTok{rng }\OperatorTok{=}\NormalTok{ np.random.default\_rng(}\DecValTok{0}\NormalTok{)}

\CommentTok{\#\# Write your data generating code here to replace the code below}
\NormalTok{n, p }\OperatorTok{=} \DecValTok{100}\NormalTok{, }\DecValTok{100}
\NormalTok{X }\OperatorTok{=}\NormalTok{ rng.normal(size}\OperatorTok{=}\NormalTok{(n, p))}
\NormalTok{y }\OperatorTok{=}\NormalTok{ rng.normal(size}\OperatorTok{=}\NormalTok{n)}
\end{Highlighting}
\end{Shaded}

I will test the dataset using the following setting:

\begin{enumerate}
\def\labelenumi{\arabic{enumi}.}
\tightlist
\item
  Split the dataset into training and test set with
  \texttt{test\_size=0.3}, \texttt{random\_state=0}.
\item
  Use \texttt{cv=5} and \texttt{neg\_mean\_squared\_error} to pick
  \(\lambda\).
\item
  The range for \texttt{alpha} is
  \texttt{alphas\ =\ np.logspace(-4,\ 3,\ 50)}.
\item
  Train the model on the training set, and test the model on the test
  set.
\item
  For performance I only look at \texttt{r2\_score}.
\end{enumerate}

\end{example}

\begin{solution}
\textbf{(a) Bias--Variance Tradeoff}

Ridge and lasso differ in how they regularize the model:

\begin{itemize}
\tightlist
\item
  Ridge uses an \(\ell_2\) penalty and shrinks all coefficients
  continuously.
\item
  Lasso uses an \(\ell_1\) penalty and can set some coefficients exactly
  to zero.
\end{itemize}

Bias

\begin{itemize}
\tightlist
\item
  Lasso may have higher bias when the true signal is dense (many small
  nonzero coefficients), since it tends to eliminate weaker predictors.
\item
  Ridge generally has lower bias in dense settings, as it retains all
  predictors.
\end{itemize}

Variance

\begin{itemize}
\tightlist
\item
  Ridge tends to have lower variance, especially when predictors are
  highly correlated, since it stabilizes all coefficients.
\item
  Lasso may have higher variance, because variable selection introduces
  instability (small changes in data can change which variables are
  selected).
\end{itemize}

Summary

\begin{itemize}
\tightlist
\item
  Ridge: lower variance, better for dense signals
\item
  Lasso: higher bias in dense settings, better for sparse signals
\end{itemize}

\textbf{(b) When Ridge Fails}

Ridge performs poorly when the true model is sparse:

\begin{itemize}
\tightlist
\item
  Only a small number of predictors are relevant
\item
  Many predictors are pure noise
\end{itemize}

In this case:

\begin{itemize}
\tightlist
\item
  Ridge keeps all variables with small coefficients
\item
  Noise variables remain and increase prediction error
\end{itemize}

In contrast:

\begin{itemize}
\tightlist
\item
  Lasso sets many coefficients to zero
\item
  Removes noise variables and improves prediction
\end{itemize}

Therefore, ridge fails when sparsity is important, while lasso succeeds
by performing variable selection.

\textbf{(c) Code example}

\begin{Shaded}
\begin{Highlighting}[]
\NormalTok{rng }\OperatorTok{=}\NormalTok{ np.random.default\_rng(}\DecValTok{0}\NormalTok{)}

\NormalTok{n, p }\OperatorTok{=} \DecValTok{120}\NormalTok{, }\DecValTok{100}
\NormalTok{X }\OperatorTok{=}\NormalTok{ rng.normal(size}\OperatorTok{=}\NormalTok{(n, p))}

\NormalTok{beta }\OperatorTok{=}\NormalTok{ np.zeros(p)}
\NormalTok{beta[:}\DecValTok{5}\NormalTok{] }\OperatorTok{=}\NormalTok{ [}\DecValTok{10}\NormalTok{, }\DecValTok{8}\NormalTok{, }\DecValTok{6}\NormalTok{, }\DecValTok{4}\NormalTok{, }\DecValTok{2}\NormalTok{]}

\NormalTok{y }\OperatorTok{=}\NormalTok{ X }\OperatorTok{@}\NormalTok{ beta }\OperatorTok{+}\NormalTok{ rng.normal(size}\OperatorTok{=}\NormalTok{n)}
\end{Highlighting}
\end{Shaded}

\begin{Shaded}
\begin{Highlighting}[]
\ImportTok{import}\NormalTok{ numpy }\ImportTok{as}\NormalTok{ np}

\NormalTok{rng }\OperatorTok{=}\NormalTok{ np.random.default\_rng(}\DecValTok{4}\NormalTok{)}
\NormalTok{n }\OperatorTok{=} \DecValTok{200}
\NormalTok{p }\OperatorTok{=} \DecValTok{500}  
\NormalTok{k }\OperatorTok{=} \DecValTok{50}  

\NormalTok{Z }\OperatorTok{=}\NormalTok{ rng.normal(size}\OperatorTok{=}\NormalTok{(n, k))}
\NormalTok{A }\OperatorTok{=}\NormalTok{ rng.normal(size}\OperatorTok{=}\NormalTok{(k, p))}
\NormalTok{X }\OperatorTok{=}\NormalTok{ Z }\OperatorTok{@}\NormalTok{ A }\OperatorTok{+} \FloatTok{0.01} \OperatorTok{*}\NormalTok{ rng.normal(size}\OperatorTok{=}\NormalTok{(n, p))  }

\NormalTok{gamma }\OperatorTok{=}\NormalTok{ rng.normal(}\DecValTok{1}\NormalTok{, }\FloatTok{0.1}\NormalTok{, size}\OperatorTok{=}\NormalTok{k)}
\NormalTok{y }\OperatorTok{=}\NormalTok{ Z }\OperatorTok{@}\NormalTok{ gamma }\OperatorTok{+} \DecValTok{3} \OperatorTok{*}\NormalTok{ rng.normal(size}\OperatorTok{=}\NormalTok{n)}
\end{Highlighting}
\end{Shaded}

\end{solution}

islr 3,4

ISlR 8 9




\end{document}
